% !PRINTED 641
\weberpage{641}
%<<CONT

Dann geht (5) in eine andere Form über, die wir so darstellen

\weq{9}{n\xi = A + \sum_{0,\,n-2}^{\nu} E_\nu K_\nu \tau_\nu^{r_{n-2}} \tau_{\nu+1}^{r_{n-3}} \ldots \tau_{\nu+n-2}^{r_0},}

worin $E_\nu$ eine $n^{\mathrm{te}}$ Einheitswurzel ist, die durch $\varepsilon$ so ausgedrückt wird:

\weq{10}{E_\nu = \varepsilon_\nu^{r_{n-2}} \varepsilon_{\nu+1}^{r_{n-3}} \ldots \varepsilon_{\nu+n-2}^{r_0}.}

Nun ist nach der Definition von $r_\nu$ [\S.\,184, (15)]

\weq{11}{r_\nu \equiv g r_{\nu-1} \pmod{n},}

und nach (10), da der Index von $r$ nach dem Modul $n-1$ zu nehmen ist,

\[
E_{\nu-1} = \varepsilon_{\nu-1}^{r_{n-2}} \varepsilon_\nu^{r_{n-3}} \ldots \varepsilon_{\nu+n-3}^{r_0}
\]
\[
E_{\nu-1}^g = \varepsilon_\nu^{r_{n-2}} \varepsilon_{\nu+1}^{r_{n-3}} \ldots \varepsilon_{\nu-1}^{r_0},
\]

also

\weq{12}{E_\nu = E_{\nu-1}^g = E_0^{r_\nu}.}

Sind also die $\varepsilon$ irgend wie bestimmt, so ergiebt sich

\weq{13}{E_0 = \varepsilon_0^{r_{n-2}} \varepsilon_1^{r_{n-3}} \ldots \varepsilon_{n-2}^{r_0},}

und dadurch sind nach (12) die übrigen $E_\nu$ vollkommen bestimmt. Also ergeben sich in der That nur $n$ Werthe aus (6), die man z.\,B. dadurch erhalten kann, dass man \gesperrt{einem} der Radicale $\tau_\nu$ seine verschiedenen Werthe beilegt. (Man vergleiche hiermit die \gesperrt{Cayley}'sche Auflösung der cubischen Gleichungen \S.\,36.)

Will man die Grössen $\xi_0, \xi_1 \ldots \xi_{n-1}$ in der Reihenfolge bestimmen, die der Bildung der cyklischen und metacyklischen Gruppe zu Grunde liegt, so muss man in (5) vor der Summation mit $\varepsilon^{-h r_\nu}$ multipliciren und findet (\S.\,133)

\weq{14}{n\xi_h = A + \sum_{0,\,n-2}^{\nu} \varepsilon^{-h r_\nu} K_\nu \tau_\nu^{r_{n-2}} \tau_{\nu+1}^{r_{n-3}} \ldots \tau_{\nu+n-2}^{r_0},}

worin eine veränderte Bestimmung der Radicale nur eine cyklische Vertauschung der $\xi_h$ bedingt.

\begin{center}
\S.\,186.

\gesperrt{Befreiung von den beschränkenden Voraussetzungen.}
\end{center}

Es wäre eine wesentliche Beschränkung dieser Untersuchungen über auflösbare Gleichungen, wenn die Voraussetzungen 1., 2. des vorigen Paragraphen aufrecht erhalten werden müssten.
